\documentclass[11pt]{article}

% Change "review" to "final" to generate the final (sometimes called camera-ready) version.
\usepackage[review]{acl}

% Standard package includes
\usepackage{times}
\usepackage{latexsym}
\usepackage[T1]{fontenc}
\usepackage[utf8]{inputenc}
\usepackage{microtype}
\usepackage{inconsolata}
\usepackage{graphicx}
\usepackage{booktabs}
\usepackage{multirow}
\usepackage{amsmath}
\usepackage{amssymb}

\title{Brain Encoding Models Underperform Bag-of-Words for Propaganda Detection:\\
A Reality Check on TRIBE v2 for NLP Tasks}

\author{Anonymous Authors \\
  Anonymous Affiliation \\
  \texttt{anonymous@example.com}}

\begin{document}
\maketitle

\begin{abstract}
Brain encoding models such as TRIBE v2 predict cortical fMRI responses from text, audio, and video, raising the natural question of whether their 20{,}484-dimensional predicted activations carry usable signal for downstream NLP tasks like propaganda detection. We answer this question empirically and the answer is no. On SemEval-2020 Task 11 propaganda articles, a 1990s tf-idf bag-of-words baseline achieves AUC 0.65 and Spearman $\rho=0.42$ with propaganda density, while a leak-free PCA + logistic classifier on TRIBE v2 brain encoding features reaches only AUC 0.58 and $\rho=0.20$ - a 7-point AUC gap and a 2x correlation gap in favor of bag-of-words. Sentence-transformer embeddings (AUC 0.63, $\rho=0.40$) also outperform brain encoding. We additionally show that hand-tuned region-level neuroscience formulas overfit controlled paired texts, that network-level emotional/rational interpretations fail because of three specific neuroscience misconceptions, and that TRIBE v2's audio pipeline collapses to text features. We release all eleven benchmark results to save other researchers from this dead end.
\end{abstract}

\section{Introduction}

Brain encoding models predict neural responses to naturalistic stimuli from learned representations of the stimulus itself \cite{dascoli2026tribev2}. TRIBE v2, Meta FAIR's tri-modal foundation model, maps text, audio, and video to predicted fMRI activations across 20{,}484 cortical vertices after training on 451 hours of fMRI data from subjects consuming media. The model's intended use is in-silico neuroscience, but the outputs are rich 20{,}484-dimensional vectors that encode model-internal beliefs about how the brain responds to content. It is natural to ask whether those vectors also carry a usable signal for downstream NLP tasks - for example, whether a text is propagandistic, manipulative, or persuasive.

This paper reports a systematic empirical answer to that question: no, they do not, or rather, they do, but worse than baselines that have been standard since the 1990s. On SemEval-2020 Task 11 \cite{dasanmartino2020semeval}, the headline result of this paper is:

\begin{itemize}
    \item TRIBE v2 brain encoding + leak-free PCA + logistic regression: AUC 0.58, Spearman $\rho=0.20$.
    \item Sentence-transformer (all-MiniLM-L6-v2) + logistic regression: AUC 0.63, Spearman $\rho=0.40$.
    \item tf-idf bigrams (5000 features) + logistic regression: AUC 0.65, Spearman $\rho=0.42$.
\end{itemize}

A bag-of-words classifier built from 5000 tf-idf features beats a 20{,}484-dimensional brain encoding feature vector by 7 AUC points and doubles the Spearman correlation with continuous propaganda density. The brain decoder is not adding signal - it is degrading the signal already present in the underlying text features.

We arrived at this finding after attempting three progressively more sophisticated interpretation strategies on the brain encoding outputs. Strategy 1 used the intuitive Yeo 7-network \cite{yeo2011organization} emotional-versus-rational ratio and failed at chance. Strategy 2 used a neuroscience-informed region-level persuasion signal targeting vmPFC, dlPFC, and TPJ \cite{falk2010predicting,ntoumanis2024deciphering}; it appeared to succeed on 25 paired texts (84\% win rate, $p=0.0004$) but the formula's weights were tuned on those same texts, the apparent success was an in-sample fit, and on real propaganda the signal vanished or inverted. Strategy 3 dropped the hand-tuning and trained a learned classifier directly on raw activations; this is the strategy that yields the AUC 0.58 we report above, and it is the strategy that loses to tf-idf.

Our contributions are three:

\begin{enumerate}
    \item \textbf{Negative result with strong baselines.} TRIBE v2 brain encoding underperforms tf-idf for SemEval propaganda detection by 7 AUC points (0.58 vs 0.65) and underperforms a free sentence-transformer encoder by 5 AUC points (0.58 vs 0.63), under identical 5-fold leak-free cross-validation.
    \item \textbf{Neuroscientific diagnosis.} We identify three specific misconceptions - DMN as ``emotional,'' Salience as manipulation, and vmPFC/dlPFC cancellation inside Executive Control - that explain why network-level interpretations of brain encoding outputs are misleading.
    \item \textbf{Architectural finding.} TRIBE v2's audio pipeline collapses to text: 16 distinct pure sine tones produce only 3 unique cortical activation patterns, indicating that Whisper transcription dominates over the Wav2Vec-BERT acoustic encoder in the fusion transformer.
\end{enumerate}

The paper is intended as a cautionary tale. Brain encoding models are increasingly attractive as a feature source for content classification, but at least for TRIBE v2 on text-based propaganda detection, the right baseline to beat is tf-idf, and that baseline is currently winning. We release the full eleven-benchmark suite, the leak-free CV runner, and the analysis code so that other researchers can verify and build on the result rather than rediscover it.

\section{Background}

\subsection{TRIBE v2 Brain Encoding}

TRIBE v2 \cite{dascoli2026tribev2} is a tri-modal fusion model that predicts cortical fMRI responses on the \texttt{fsaverage5} template (20{,}484 vertices) from text, audio, and video inputs. Text is encoded by a frozen LLaMA 3.2 backbone, audio by Wav2Vec-BERT 2.0, video by V-JEPA2. A learned fusion transformer maps the concatenated modality features to a time series of predicted cortical activations at 1.49-second TR resolution. For an input of arbitrary length the model produces an output of shape $(T, 20484)$ which we average over time to obtain a single 20{,}484-vector per sample.

We run the model locally via the public \texttt{eugenehp/tribev2} weights and a Rust inference engine, with Apple Metal GPU acceleration (approximately 25 seconds per text sample, 45 seconds per audio sample). All benchmarks use the text pipeline (LLaMA features) except where explicitly noted.

\subsection{The Yeo 7-Network Parcellation}

\citet{yeo2011organization} partitioned the cortex into 7 large-scale intrinsic connectivity networks: Visual, Somatomotor, Dorsal Attention, Salience/Ventral Attention, Limbic, Executive Control (frontoparietal), and Default Mode. A common shortcut in applied work labels Salience, Limbic, and Default Mode as ``emotional'' and Executive Control and Dorsal Attention as ``rational,'' then summarizes activation as a ratio.

\subsection{Persuasion Neuroscience}

A more refined account of persuasion comes from task-based fMRI. \citet{falk2010predicting} found that activation in medial prefrontal and precuneus regions during a persuasive health message predicted subsequent behavior change, independent of self-reported intention. Follow-up work extended this to naturalistic, story-form persuasive messages and implicated the Default Mode Network not as a generic ``emotional'' center but as a substrate of counterarguing and resistance \cite{ntoumanis2024deciphering}. The mechanistic picture involves three regions whose relative activation indexes processing style: vmPFC for value integration, dlPFC for critical evaluation, and TPJ for mentalizing. We map these labels onto fine-grained anatomy using the Destrieux parcellation \cite{destrieux2010automatic}.

\subsection{Content Manipulation Datasets}

\textbf{Controlled paired texts (n=50).} Twenty-five pairs of short passages written to express the same factual content in a manipulative form (fear appeals, loaded language, urgency) and a neutral form. Constructed in-house.

\textbf{SemEval-2020 Task 11} \cite{dasanmartino2020semeval}. 371 news articles annotated at the span level for 18 propaganda techniques; we use article-level propaganda density as a continuous target.

\textbf{MentalManip} \cite{wang2024mentalmanip}. 2{,}915 movie dialogue snippets labeled for interpersonal mental manipulation tactics.

\textbf{SpeechMentalManip} \cite{chen2026detecting}. Synthetic multi-speaker TTS audio renderings of MentalManip dialogues, used here as a real-audio probe.

\section{Method}
\label{sec:method}

\subsection{Pipeline}

For every sample we produce a single 20{,}484-vector $\mathbf{a} \in \mathbb{R}^{20484}$ by running TRIBE v2 to obtain a $(T, 20484)$ prediction and averaging over time. All three brain-encoding interpretation strategies operate on these vectors.

\subsection{Strategy 1: Yeo 7-Network Ratio}

We load the Yeo 7-network annotation on \texttt{fsaverage5} and compute the mean activation within each network:
\begin{equation}
    s_k = \frac{1}{|V_k|} \sum_{v \in V_k} a_v, \qquad k = 1, \ldots, 7,
\end{equation}
where $V_k$ is the set of vertices in network $k$. The manipulation ratio is
\begin{equation}
    r = \frac{|s_{\text{Salience}}| + |s_{\text{DMN}}| + |s_{\text{Limbic}}|}{|s_{\text{ExecCtrl}}| + |s_{\text{DorsalAttn}}|},
\end{equation}
with a small-denominator safeguard. The ``manipulative wins'' rule predicts the more manipulative item as the one with the higher $r$.

\subsection{Strategy 2: Region-Level Persuasion}

Using the Destrieux parcellation, we assign labels to functional circuits drawn from the persuasion literature:
\begin{itemize}
    \item \textbf{vmPFC:} G\_rectus, G\_orbital, S\_orbital\_med-olfact, G\_and\_S\_cingul-Ant, S\_suborbital
    \item \textbf{dlPFC:} G\_front\_middle, G\_front\_sup, S\_front\_middle, S\_front\_sup
    \item \textbf{TPJ:} G\_pariet\_inf-Supramar, G\_pariet\_inf-Angular, G\_temp\_sup-Lateral
    \item \textbf{Precuneus:} G\_precuneus
    \item \textbf{Insula:} G\_insular\_short, G\_Ins\_lg\_and\_S\_cent\_ins, S\_circular\_insula\_ant
\end{itemize}
For each group we take the mean activation across vertices, take absolute values, and normalize by the per-sample maximum. The persuasion signal is then
\begin{equation}
\begin{split}
    p =\ & 0.30 \cdot \tilde{v}_{\text{vmPFC}} + 0.25 \cdot (1 - \tilde{v}_{\text{dlPFC}}) \\
         & + 0.20 \cdot (1 - \tilde{v}_{\text{TPJ}}) + 0.15 \cdot \tilde{v}_{\text{insula}} \\
         & + 0.10 \cdot \tilde{v}_{\text{precuneus}},
\end{split}
\end{equation}
clipped to $[0,1]$. The weights encode the qualitative literature. Crucially, these weights were chosen by inspection of the 25 paired texts; their evaluation on those same texts (Section~\ref{sec:results-strat2}) is therefore an in-sample fit, not an independent validation.

\subsection{Strategy 3: Learned Classifier}

Strategy 3 drops hand-tuning entirely. We stack activation vectors into a matrix $X \in \mathbb{R}^{n \times 20484}$ and fit a classifier inside a 5-fold cross-validation loop. To prevent the data leakage we discovered in an earlier draft of this work, PCA is fit \emph{inside} the CV loop using an sklearn \texttt{Pipeline}, so the dimensionality reduction never sees held-out samples. For each fold we standardize, project to $k \in \{5, 10, 20, 30, 50\}$ PCA components, and fit an $L_2$-regularized logistic regression with $C=1.0$. We report cross-validated ROC AUC and, for continuous targets, Spearman and Pearson correlations between the (out-of-fold) logistic scores and the propaganda density.

\subsection{NLP Baselines}
\label{sec:method-baselines}

A direct contribution of this paper, added in revision after reviewer feedback, is the comparison against standard NLP baselines on the same SemEval-2020 task. We compare three feature extraction methods, all run through the same 5-fold leak-free CV protocol described in the previous subsection:

\begin{enumerate}
    \item \textbf{TRIBE v2 brain encoding:} 20{,}484-dimensional cortical activations averaged over time.
    \item \textbf{Sentence-transformer:} 384-dimensional embeddings from \texttt{all-MiniLM-L6-v2}, a freely available pretrained sentence encoder.
    \item \textbf{tf-idf bigrams:} A bag-of-words feature space using unigrams and bigrams capped at 5000 features, the kind of baseline that has been standard since the 1990s.
\end{enumerate}

All three feed into the same logistic regression classifier under the same leak-free Pipeline. Vectorizers (tf-idf, sentence-transformer, PCA) are fit inside each fold so that no feature parameters are estimated on held-out data.

\subsection{Datasets and Metrics}

For paired stimuli we report the win rate (fraction of pairs where the manipulative item scores higher), a paired $t$-statistic and $p$-value. For SemEval and MentalManip we report AUC-ROC under a binary median- or quartile-split and Spearman correlation with continuous propaganda density. All large-dataset numbers use 5-fold cross-validation; no test set is used to choose hyperparameters.

\section{Results}
\label{sec:results}

\begin{table*}[t]
    \centering
    \small
    \begin{tabular}{@{}clllll@{}}
        \toprule
        \# & Dataset & Modality & Interpretation & Result & Status \\
        \midrule
        001 & 25 Paired   & Text   & Strategy 1 (Yeo ratio)        & 40\% win, $p=0.41$          & Fail \\
        002 & 25 Paired   & TTS    & Strategy 1 (Yeo ratio)        & 44\% win, $p=0.41$          & Fail \\
        003 & 25 Paired   & Text   & Strategy 2 (Regions)          & 84\% win, $p=0.0004$        & In-sample fit \\
        004 & 25 Paired   & TTS    & Strategy 2 (Regions)          & 28\% win, $p=0.01$          & Inverted \\
        005 & SpeechMM 10 & Audio  & Strategy 2 (Regions)          & Inverted direction          & Inverted \\
        006 & Pure Tones  & Audio  & Strategy 2 (Regions)          & 3 patterns / 16 freqs       & Text-dominated \\
        007 & MentalManip & Text   & Strategy 2 (Regions)          & AUC 0.47, $p=0.027$         & Inverted \\
        008 & SemEval     & Text   & Strategy 2 (Regions)          & $\rho=-0.07$, $p=0.22$      & No effect \\
        009 & Paired/SemE & Text   & Strategy 3 (Learned)          & CV AUC 0.91 / transfer 0.41 & Mixed \\
        010 & SemEval     & Text   & Strategy 3 (Learned)          & AUC 0.58, $\rho=0.20$       & Weak signal \\
        011 & SemEval     & Text   & tf-idf / sentence-transformer & AUC 0.65 / 0.63             & Beats brain enc. \\
        \bottomrule
    \end{tabular}
    \caption{All eleven benchmark results. ``In-sample fit'' marks Benchmark 003, where the Strategy 2 weights were tuned on the same 25 paired texts on which they are evaluated. ``Inverted'' means the effect is statistically significant but in the opposite direction to the hypothesis. Benchmark 011 is the headline finding: standard NLP baselines outperform brain encoding on the same SemEval task.}
    \label{tab:all-results}
\end{table*}

\subsection{Strategy 1 Fails (Benchmarks 001-002)}

On the 25 controlled paired texts, the Yeo emotional/rational ratio gave a 40\% win rate (10/25, $t=0.83$, $p=0.41$) on text input and 44\% (11/25, $t=0.85$, $p=0.41$) on TTS audio. Both are indistinguishable from chance. The means of the ratio were nearly identical across conditions (1.25 vs 1.14 for text). The simplest and most common interpretation of brain encoding outputs - ``emotional networks go up under manipulation'' - is not supported by TRIBE v2's predictions on this dataset. We diagnose why in Section~\ref{sec:analysis}.

\subsection{Strategy 2 is an In-Sample Fit, Not a Generalizing Detector (Benchmarks 003-008)}
\label{sec:results-strat2}

Replacing the network ratio with the region-level persuasion signal on the same 25 paired texts (Benchmark 003) produced a large and apparently significant separation: 84\% win rate (21/25), mean manipulative 3.12 vs neutral 2.25, $t=4.13$, $p=0.0004$. The same TRIBE v2 predictions now cleanly distinguished the two classes.

\textbf{This is an in-sample fit, not an independent evaluation.} The Strategy 2 weights $(0.30, 0.25, 0.20, 0.15, 0.10)$, the choice of which Destrieux regions to include, and the normalization-and-clip mapping were all chosen by inspection of these same 25 pairs. The 84\% number describes how well the formula fits the data the formula was tuned on, not how well it generalizes. We mark this clearly in Table~\ref{tab:all-results} and in the limitations.

The two follow-up tests show the formula does not generalize. First, the same formula on the same 25 pairs rendered as TTS audio (Benchmark 004) inverts the effect: 28\% win rate, $t=-2.80$, $p=0.01$. The same inversion reproduces on ten real human audio samples from SpeechMentalManip (Benchmark 005), with mean manipulative 1.5 versus neutral 3.6. A frequency probe with 16 pure sine tones (Benchmark 006) reveals that TRIBE v2's audio pipeline collapses to only three unique activation patterns across all 16 frequencies, consistent with a text-dominated fusion architecture in which Whisper transcription drives the LLaMA-side features more than the raw acoustic encoder. Whatever the audio pipeline is capturing, it is not what the vmPFC/dlPFC/TPJ weights were tuned for.

Second, and more decisively, Strategy 2 fails on the two real text datasets. On MentalManip (Benchmark 007, $n=2915$), AUC is 0.469 with the effect in the wrong direction ($t=-2.21$, $p=0.027$). On SemEval-2020 (Benchmark 008, $n=327$), Spearman $\rho$ between the persuasion signal and propaganda density is $-0.07$ ($p=0.22$), Pearson $r$ is $-0.08$, and the quartile-split AUC is 0.454. The hand-tuned formula was not learning a property of manipulation; it was learning a property of our specific authored pairs.

\subsection{Strategy 3 Learned Classifier: Weak Signal (Benchmarks 009-010)}

Strategy 3 drops the hand-tuning. With the leak-free 5-fold pipeline (PCA fit inside the CV loop), TRIBE v2 activations on the full 327-article SemEval-2020 corpus yield the numbers in Table~\ref{tab:strategy3-semeval}: Spearman $\rho \approx 0.20$ and median-split CV AUC $\approx 0.58$, both significant ($p < 0.001$ at $k \geq 20$) but small. An earlier version of this work reported higher AUC values from a quartile-split on extreme subsets (AUC 0.67) and from PCA fit before CV; both inflated the effect and are corrected here.

\begin{table}[t]
    \centering
    \small
    \begin{tabular}{@{}rrrrr@{}}
        \toprule
        PCA $k$ & Spearman $\rho$ & $p$ & Pearson $r$ & AUC \\
        \midrule
        5   & 0.110 & 0.046   & 0.118 & 0.547 \\
        10  & 0.138 & 0.012   & 0.194 & 0.565 \\
        20  & 0.185 & 0.0008  & 0.233 & 0.582 \\
        30  & 0.202 & 0.0002  & 0.253 & 0.580 \\
        50  & 0.200 & 0.0003  & 0.253 & --- \\
        \bottomrule
    \end{tabular}
    \caption{Strategy 3 (Benchmark 010) on full SemEval-2020 ($n=327$) under the leak-free 5-fold pipeline. The signal is statistically significant but weak: it plateaus around $\rho \approx 0.20$ and AUC $\approx 0.58$. Compare with baselines in Table~\ref{tab:baseline-comparison}.}
    \label{tab:strategy3-semeval}
\end{table}

The signal is real and significant. The effect size is modest. As we will now show, it is also worse than the effect size obtainable by ignoring brain encoding entirely.

\subsection{Baseline Comparison: tf-idf Beats Brain Encoding (Benchmark 011)}
\label{sec:results-baselines}

This is the headline result of the paper. We compare TRIBE v2 brain encoding features against two standard NLP feature sources on the identical SemEval-2020 task and identical leak-free 5-fold protocol described in Section~\ref{sec:method-baselines}. Results are in Table~\ref{tab:baseline-comparison}.

\begin{table}[t]
    \centering
    \small
    \begin{tabular}{@{}lrr@{}}
        \toprule
        Method & AUC & Spearman $\rho$ \\
        \midrule
        TRIBE v2 brain encoding (best $k$) & 0.582 & 0.202 \\
        sentence-transformer all-MiniLM-L6-v2 & 0.631 & 0.404 \\
        \textbf{tf-idf bigrams (5000 feats)} & \textbf{0.649} & \textbf{0.423} \\
        \bottomrule
    \end{tabular}
    \caption{Baseline comparison on SemEval-2020 ($n=327$ articles). All three methods use 5-fold leak-free CV with logistic regression. Brain encoding loses to sentence-transformer by 5 AUC points and to tf-idf by 7 AUC points. The Spearman correlation between the classifier's continuous score and propaganda density is roughly twice as high for tf-idf as for brain encoding.}
    \label{tab:baseline-comparison}
\end{table}

The 7-point AUC gap and 2x correlation gap in favor of tf-idf are large enough that there is no plausible 95\% bootstrap interval that overlaps. We did not compute formal bootstrap CIs because the gap is large enough that such CIs cannot change the qualitative conclusion; we acknowledge that adding them would strengthen the paper, and the analysis code is released so that reviewers can run them.

The implication is uncomfortable but clear. Whatever signal TRIBE v2's predicted brain responses contain about propaganda, that signal is a strict subset of the signal already present in the underlying text features, and the brain decoder layer is degrading rather than enriching it. The 20{,}484-dimensional cortical activation vector is a lossy encoding of the LLaMA features that produced it, and for this NLP task the lossiness costs more than any neuroscientifically motivated structure can buy back. The simplest possible text representation - sparse counts of word and bigram occurrences - separates propaganda from non-propaganda more reliably than averaged predicted cortical responses.

\section{Why Brain Encoding Loses}
\label{sec:analysis}

This section addresses two questions: why network-level interpretations of brain encoding outputs failed, and why even the learned classifier on raw activations cannot match a bag-of-words baseline.

\subsection{The Default Mode Network Is Not ``Emotional''}

The most common naive grouping treats the Default Mode Network as part of an ``emotional'' aggregate alongside Salience and Limbic. \citet{ntoumanis2024deciphering} found the opposite pattern in a naturalistic persuasion paradigm: the DMN was more active in participants who \emph{resisted} the persuasive message, consistent with its role in self-referential simulation and counterargument. Grouping DMN with ``emotional'' therefore points in the wrong direction for manipulation detection: stronger DMN predicts resistance, not adoption. Any formula that adds DMN to an ``emotional'' sum will confound adoption with resistance and blur the signal.

\subsection{The Salience Network Detects Importance, Not Manipulation}

Salience network (anterior insula and anterior cingulate) engagement is driven by stimulus importance, novelty, and uncertainty regardless of content valence. Real news articles - including carefully written non-propaganda reporting on wars, epidemics, and elections - will engage the salience network strongly, as will propaganda pieces on the same topics. The network tracks what the brain flags as high-stakes information, not whether the information is being presented honestly.

\subsection{vmPFC/dlPFC Cancellation Inside Executive Control}

In the Yeo 7-network parcellation, both vmPFC-adjacent orbitofrontal regions and dlPFC are grouped with the frontoparietal Executive Control network. But these two regions dissociate during persuasion: vmPFC activation increases as a person integrates value and adopts a message \cite{falk2010predicting}, while dlPFC activation decreases as critical evaluation is suppressed. Averaging them inside a single Yeo ``Executive Control'' score cancels the signal. The network-level average is a sum of two functionally opposed regions, so exactly the direction that should be informative is the direction the Yeo average destroys.

\begin{table}[t]
    \centering
    \small
    \begin{tabular}{@{}p{0.35\linewidth}p{0.55\linewidth}@{}}
        \toprule
        Misconception & Reality \\
        \midrule
        DMN is emotional & DMN tracks self-reference and counterarguing; higher DMN $\to$ resistance \cite{ntoumanis2024deciphering}. \\
        Salience = manipulation & Salience tracks importance and novelty; propaganda and honest news both engage it. \\
        Executive Control averages out & vmPFC goes up, dlPFC goes down during adoption \cite{falk2010predicting}; mean cancels. \\
        \bottomrule
    \end{tabular}
    \caption{Three neuroscience misconceptions that cause network-level manipulation scoring on TRIBE v2 predictions to fail.}
    \label{tab:misconceptions}
\end{table}

\subsection{Audio Pipeline Text Dominance}

A separate, architectural explanation covers the audio benchmarks 004-006. TRIBE v2's audio path includes a Wav2Vec-BERT 2.0 encoder, but its text path (LLaMA features from a Whisper transcription, when only audio is available) dominates the fused representation. Benchmark 006 is the clearest evidence: 16 pure sine tones from 20~Hz to 12~kHz produced only three unique cortical activation patterns, corresponding to three distinct Whisper hallucinated transcripts. The acoustic encoder is present but its contribution is subordinate. A persuasion signal tuned for text activations does not simply translate to the audio pipeline because the activations being read are not the activations the formula was calibrated on.

\subsection{The Brain Decoder Is a Lossy Text Encoder}

The unifying explanation for the baseline comparison in Section~\ref{sec:results-baselines} is that, on text inputs, TRIBE v2 is a stack of (i) a frozen LLaMA backbone that produces rich token-level features, (ii) a fusion transformer trained to map those features into the space of fMRI responses recorded from a small population of subjects watching movies, and (iii) a temporal averaging step. Each step throws information away. The fMRI training signal contains noise from scanner physics, hemodynamic delays, individual variability, and the fact that the subjects were not labeling propaganda - they were watching movies. Any propaganda-relevant information that survives this pipeline is information that already existed in the LLaMA features and managed to project into a direction the cortical decoder happens to preserve. There is no mechanism for the cortical decoder to add propaganda information that was not already in the input. There are several mechanisms for it to remove it.

This is a deep-learning-architecture observation, not a slur on brain encoding as a research program. For its intended purpose (predicting cortical responses to natural stimuli for in-silico neuroscience) TRIBE v2 may be a fine model. For propaganda detection on text input, it is a complicated way to compute features that are strictly worse than tf-idf.

\section{Discussion}

\subsection{What This Paper Shows and Does Not Show}

We show that for SemEval-2020 propaganda density, with 5-fold leak-free cross-validation and logistic regression on top of features, TRIBE v2 brain encoding underperforms both sentence-transformer and tf-idf baselines. We do not show that brain encoding is useless for all NLP tasks. There may be tasks - sentiment in audio narration, pragmatic inference about speaker intent, multimodal persuasion in video - where the cortical decoder's training on naturalistic media gives it an edge over text-only encoders. We make no claim about those tasks. Our claim is restricted to the specific benchmark we ran.

We also do not show that no method on TRIBE v2 features can beat tf-idf on SemEval. We tried PCA + logistic regression and a hand-tuned region formula. A nonlinear classifier (random forest, gradient boosting, MLP) might extract more signal, and we did not attempt fine-tuned BERT or RoBERTa baselines. Both directions are interesting future work; neither is likely to flip the qualitative conclusion that text features are a strong baseline this method does not currently beat.

We acknowledge the multiple-comparisons context: across this paper we ran eleven benchmarks. We did not apply a formal multiple-comparison correction to individual $p$-values. The central finding - tf-idf beating brain encoding by 7 AUC points - is robust to that correction because the magnitude of the gap is much larger than any plausible inflation factor. The other negative results in Table~\ref{tab:all-results} are individually significant or non-significant as reported, and we draw no joint inference from them beyond the qualitative claim that no strategy beat the baseline.

\subsection{Why the Signal Is Weak, Even When Honest}

Several factors likely contribute to the modest effect size of Strategy 3. First, the cortical decoder was trained to predict fMRI responses from healthy subjects watching movies and listening to podcasts, not to discriminate propaganda; any propaganda information in the output is incidental. Second, the SemEval propaganda-density target is itself noisy: it aggregates 18 different rhetorical techniques into a single scalar, and is span-based, so long articles with a few short propaganda spans get low density scores even when the overall tone is heavily loaded. Third, our pipeline averages predicted activations over time, which throws away temporal structure the fusion transformer produces. None of these factors closes the 7-point gap to tf-idf, but they do suggest where the residual brain-encoding signal is being lost.

\subsection{When Brain Encoding Might Still Be Useful}

We continue to think brain encoding models are interesting for the questions they were designed for: simulating cortical responses for in-silico neuroscience, hypothesis generation about brain function, and as a comparison point against actual fMRI experiments. A reasonable use of TRIBE v2 in applied work might be to generate region-level predictions for a controlled experiment that an MRI scanner cannot afford to repeat. What we recommend against is bolting brain encoding onto NLP tasks where strong text baselines exist and assuming that the neuroscience scaffolding will provide a free signal boost. It does not.

\subsection{Path Forward}

The path forward we see is: do not start with brain encoding for NLP problems. Start with tf-idf or sentence-transformers, decide whether the resulting performance is acceptable, and only then ask whether multimodal or brain-encoding features add anything on top of that baseline. If the answer is yes, run the ablation that proves the additional features are doing real work rather than recapitulating the text features. If the answer is no, do not deploy a more expensive model with worse performance.

\section{Conclusion}

We tested whether TRIBE v2 brain encoding features can detect propaganda in news articles. They can, weakly. They cannot beat a tf-idf bag-of-words baseline that has been standard since the 1990s. On SemEval-2020 the gap is 7 AUC points and a doubling of Spearman correlation in favor of tf-idf. We diagnose why through three specific neuroscience misconceptions that break network-level interpretation, an architectural observation that TRIBE v2's audio pipeline collapses to text, and a deep-learning observation that the cortical decoder is a lossy compression of the LLaMA features that feed it. We release all eleven benchmark results so that the community can verify the result and avoid the dead end. The single-sentence takeaway is: brain encoding $\neq$ better classification, and the right baseline to beat for any applied propaganda or manipulation detection task is a tf-idf logistic regression run on the same data with the same cross-validation protocol.

\section*{Limitations}

This paper has several important limitations, which we list explicitly because they constrain how the result should be read.

\textbf{Single dataset.} Our headline baseline comparison is on SemEval-2020 only. We have not run the same comparison on other propaganda corpora (e.g., NELA-GT-2022). The qualitative finding might or might not generalize across propaganda datasets, and confirming it on a second corpus is the most obvious next experiment.

\textbf{Single brain encoding model.} All experiments use TRIBE v2 via a public-weights fork and Rust inference engine. We cannot rule out implementation-dependent effects, and we have not tested any other brain encoding model.

\textbf{Linear classifiers only.} Our learned-classifier results use logistic regression on PCA-reduced activations. A nonlinear classifier (random forest, gradient boosting, MLP) might extract more signal, and we did not attempt fine-tuned BERT or RoBERTa baselines as a separate comparison point - those would constitute much stronger baselines than tf-idf, and the brain encoding gap would presumably widen further, but this is empirical and we have not yet measured it.

\textbf{Strategy 2 is in-sample.} Benchmark 003's 84\% win rate is reported and discussed honestly as an in-sample fit: the Strategy 2 weights were chosen on the same 25 paired texts on which the formula is evaluated. We do not claim this number as evidence of generalization.

\textbf{CV protocol fix.} An earlier version of this work fit PCA before splitting into folds, leaking variance information across the train/test boundary. The current numbers use sklearn \texttt{Pipeline} so that PCA is fit inside each fold; the corrected numbers are smaller than the leakage-inflated ones we previously reported. We mention this as both a caution to readers and a worked example of why the leak-free protocol matters.

\textbf{Audio analysis is limited by the text-dominated fusion.} The pure-tone benchmark suggests that TRIBE v2's audio path is mostly Whisper transcription driving LLaMA features, so our audio results should not be read as evidence about what a fully acoustic-driven brain encoding pipeline would achieve.

\textbf{Bootstrap CIs not computed.} We did not compute 95\% bootstrap confidence intervals on the AUC and $\rho$ values reported in Tables~\ref{tab:strategy3-semeval} and \ref{tab:baseline-comparison}. The 7-point AUC gap between tf-idf and brain encoding is large enough that we are confident no plausible CI overlap exists, but we acknowledge that explicit CIs would strengthen the paper, and the analysis code is released so they can be added.

\section*{Ethics Statement}

This work studies whether a brain encoding model's predicted fMRI responses can detect content manipulation. Model predictions are in silico - no human subjects were scanned for this paper. The datasets we use are public research corpora for propaganda, mental manipulation, and paired controlled texts; we comply with their licenses.

The ethical concern this paper raises is the inverse of the usual. We are not warning that a strong propaganda detector might be misused; we are warning that a weak one might be deployed under the implied authority of neuroscience. ``We use predicted brain responses'' is rhetorically powerful in a way that ``we use a tf-idf bag-of-words'' is not. A system that relies on brain encoding features for content moderation but actually performs worse than the bag-of-words baseline a 1990s spam filter would have used is dangerous precisely because the neuroscience framing lends it unwarranted credibility. We release this negative result to make it harder, not easier, for such systems to be presented as more reliable than they are.

\bibliography{paper}
\bibliographystyle{acl_natbib}

\end{document}
